% Sections 5.1 to 5.8, from lectures/L05/appendix/a_theory.md.

\section{Potenser}
\label{c5:sec:potenser}
En \textbf{potens} är upprepad multiplikation av ett tal med sig självt:
\[
  a^n = \underbrace{a \times a \times \cdots \times a}_{n \text{ faktorer}}
\]
Här kallas $a$ \textbf{bas} och $n$ \textbf{exponent}.

\begin{exempel}
\[
  2^4 = 2 \times 2 \times 2 \times 2 = 16, \qquad 10^3 = 1\,000
\]
\end{exempel}

\section{Potensregler}
\label{c5:sec:potensregler}
Följande sju regler gäller för alla $a \neq 0$ och $b \neq 0$.

\begingroup
\renewcommand{\arraystretch}{1.7}
\begin{narrowtable}{@{}cll@{}}
  \thead{Nr} & \thead{Regel} & \thead{Formel} \\
  \midrule
  1 & Produktregeln & $a^m \cdot a^n = a^{m+n}$ \\
  2 & Kvotregeln & $\dfrac{a^m}{a^n} = a^{m-n}$ \\
  3 & Potens av potens & $(a^m)^n = a^{m \cdot n}$ \\
  4 & Potens av produkt & $(a \cdot b)^n = a^n \cdot b^n$ \\
  5 & Potens av kvot & $\left(\dfrac{a}{b}\right)^n = \dfrac{a^n}{b^n}$ \\
  6 & Nollexponent & $a^0 = 1$ \\
  7 & Negativ exponent & $a^{-n} = \dfrac{1}{a^n}$ \\
\end{narrowtable}
\endgroup

\begin{exempel}[Regel 1]
\[
  10^3 \cdot 10^4 = 10^{3+4} = 10^7
\]
\end{exempel}

\begin{exempel}[Regel 7]
\[
  10^{-3} = \frac{1}{10^3} = \frac{1}{1000} = 0{,}001
\]
\end{exempel}

\section{Rötter}
\label{c5:sec:rotter}
\textbf{Kvadratroten} av $a$ ($a \geq 0$) är det icke-negativa tal vars kvadrat är $a$:
\[
  \sqrt{a} = a^{1/2} \quad \text{eftersom} \quad \left(\sqrt{a}\right)^2 = a
\]
\textbf{$n$:te roten} av $a$ är det tal vars $n$:te potens är $a$:
\[
  \sqrt[n]{a} = a^{1/n}
\]

\subsection{Bråkexponenter}
\label{c5:sec:brakexponenter}
Rötter och potenser kombineras med bråkexponenter:
\[
  a^{m/n} = \sqrt[n]{a^m} = \left(\sqrt[n]{a}\right)^m
\]

\begin{exempel}
\[
  8^{2/3} = \left(\sqrt[3]{8}\right)^2 = 2^2 = 4
\]
\end{exempel}

\subsection{Räkneregler för rötter}
\label{c5:sec:rotregler}
\[
  \sqrt{a \cdot b} = \sqrt{a} \cdot \sqrt{b}, \qquad \sqrt{\frac{a}{b}} = \frac{\sqrt{a}}{\sqrt{b}}
\]

\begin{aside}
\note[Obs] $\sqrt{a + b} \neq \sqrt{a} + \sqrt{b}$. Det är ett vanligt misstag!
\end{aside}

\section{Standardform (vetenskaplig notation)}
\label{c5:sec:standardform}
I elektroteknik förekommer extremt stora och extremt små tal, till exempel
$R = 1\,000\,000\,\Omega$ och $C = 0{,}000\,000\,001\,\text{F}$.

\textbf{Standardform} skriver talet som $a \times 10^n$, där $1 \leq a < 10$:
\[
  1\,000\,000 = 1{,}0 \times 10^6, \qquad 0{,}000\,000\,001 = 1{,}0 \times 10^{-9}
\]

\subsection{SI-prefix}
\label{c5:sec:siprefix}

\begin{narrowtable}{@{}llc@{}}
  \thead{Prefix} & \thead{Symbol} & \thead{Faktor} \\
  \midrule
  Giga & G & $10^9$ \\
  Mega & M & $10^6$ \\
  Kilo & k & $10^3$ \\
  Milli & m & $10^{-3}$ \\
  Mikro & $\mu$ & $10^{-6}$ \\
  Nano & n & $10^{-9}$ \\
  Piko & p & $10^{-12}$ \\
\end{narrowtable}

\begin{exempel}
$4{,}7\,\text{k}\Omega = 4{,}7 \times 10^3\,\Omega = 4\,700\,\Omega$
\end{exempel}

\begin{exempel}
$33\,\text{nF} = 33 \times 10^{-9}\,\text{F} = 3{,}3 \times 10^{-8}\,\text{F}$
\end{exempel}

\section{Räkning med standardform}
\label{c5:sec:rakning}
\textbf{Multiplikation:}
\[
  (3 \times 10^4) \times (2 \times 10^3) = 6 \times 10^7
\]
\textbf{Division:}
\[
  \frac{6 \times 10^6}{2 \times 10^2} = 3 \times 10^4
\]
\textbf{Addition} (termerna måste ha samma exponent):
\[
  3{,}5 \times 10^3 + 2{,}0 \times 10^3 = 5{,}5 \times 10^3
\]

\section{Värdesiffror och avrundning}
\label{c5:sec:vardesiffror}
\textbf{Värdesiffror} är de siffror som bär meningsfull information i ett mätvärde.
\begin{itemize}
  \item $4\,700\,\Omega$ kan ha 2, 3 eller 4 värdesiffror; det går inte att avgöra utan
    sammanhang.
  \item $4{,}70 \times 10^3\,\Omega$ har tydligt \textbf{3 värdesiffror}.
\end{itemize}
Resultatet av en beräkning bör normalt anges med lika många värdesiffror som det minst exakta
indata.

\textbf{Avrundningsregel:} Om siffran efter avrundningspositionen är $5$ eller större avrundas
uppåt, annars nedåt.

\section{Tillämpning: effektberäkning}
\label{c5:sec:effekt}
Effekten som utvecklas i ett motstånd $R$ vid spänningen $U$ är
\[
  P = \frac{U^2}{R}.
\]

\begin{exempel}
Med $U = 6\,\text{V}$ och $R = 50\,\Omega$ blir
\[
  P = \frac{6^2}{50} = \frac{36}{50} = 0{,}72\,\text{W}.
\]
\end{exempel}

Vid en given ström $I$ är effekten
\[
  P = RI^2.
\]

\begin{exempel}
Med $I = 20\,\text{mA} = 20 \times 10^{-3}\,\text{A}$ och $R = 1\,\text{k}\Omega = 10^3\,\Omega$
blir
\[
  P = 10^3 \times (20 \times 10^{-3})^2 = 10^3 \times 400 \times 10^{-6} = 0{,}4\,\text{W}.
\]
\end{exempel}

\section{Binära och hexadecimala tal}
\label{c5:sec:binart}
Vårt vanliga talsystem är ett \textbf{positionssystem} med basen $10$: en siffras värde beror på
var i talet den står, och varje position motsvarar en tiopotens. Positionen längst till höger har
värdet $10^0 = 1$:
\[
  4\,703 = 4 \cdot 10^3 + 7 \cdot 10^2 + 0 \cdot 10^1 + 3 \cdot 10^0
\]
Samma idé fungerar med vilken bas som helst. Inom digitalteknik och programmering används, utöver
det decimala talsystemet, även det \textbf{binära} och det \textbf{hexadecimala}. Tabellen visar
talet $45$ i alla tre:

\begin{narrowtable}{@{}lcll@{\hspace{2em}}l@{}}
  \thead{Talsystem} & \thead{Bas} & \thead{Siffror} & \thead{Skrivsätt} & \thead{I C} \\
  \midrule
  Decimalt & $10$ & 0--9 & $45$ & \code{45} \\
  Binärt & $2$ & 0 och 1 & $101101_2$ & \code{0b101101} \\
  Hexadecimalt & $16$ & 0--9 och A--F & $\mathrm{2D}_{16}$ & \code{0x2D} \\
\end{narrowtable}

De hexadecimala siffrorna A--F betyder $10$--$15$. Den lilla siffran efter talet anger basen;
saknas den är talet decimalt.

En binär siffra kallas \textbf{bit}, och åtta bitar kallas en \textbf{byte}. Med $n$ bitar kan
$2^n$ olika värden lagras, från $0$ till $2^n - 1$. En byte rymmer alltså $2^8 = 256$ olika
värden, $0$--$255$.

\subsection{Från binärt eller hexadecimalt till decimalt}
\label{c5:sec:tilldecimalt}
Multiplicera varje siffra med positionens värde och summera. Positionerna numreras från höger med
början på $0$, så position $n$ har värdet $2^n$ i ett binärtal och $16^n$ i ett hexadecimalt tal:

\begin{narrowtable}{@{}lcccccccc@{}}
  \thead{Position $n$} & $7$ & $6$ & $5$ & $4$ & $3$ & $2$ & $1$ & $0$ \\
  \midrule
  $2^n$ & $128$ & $64$ & $32$ & $16$ & $8$ & $4$ & $2$ & $1$ \\
\end{narrowtable}

För hexadecimala tal räcker oftast $16^0 = 1$, $16^1 = 16$, $16^2 = 256$ och $16^3 = 4\,096$.

\begin{exempel}
\[
  1011_2 = 1 \cdot 2^3 + 0 \cdot 2^2 + 1 \cdot 2^1 + 1 \cdot 2^0 = 8 + 0 + 2 + 1 = 11
\]
\end{exempel}

\begin{exempel}
F betyder $15$:
\[
  \mathrm{2F}_{16} = 2 \cdot 16^1 + 15 \cdot 16^0 = 32 + 15 = 47
\]
\end{exempel}

\subsection{Från decimalt till binärt eller hexadecimalt}
\label{c5:sec:frandecimalt}
Dividera talet med basen upprepade gånger och anteckna resten varje gång, tills kvoten blir $0$.
Resterna, lästa \textbf{nerifrån och upp}, är talet i den nya basen.

\begin{exempel}[Från decimalt till binärt]
Omvandla $13$ till binär form.

\begin{narrowtable}{@{}ccc@{}}
  \thead{Division} & \thead{Kvot} & \thead{Rest} \\
  \midrule
  $13 / 2$ & $6$ & $1$ \\
  $6 / 2$ & $3$ & $0$ \\
  $3 / 2$ & $1$ & $1$ \\
  $1 / 2$ & $0$ & $1$ \\
\end{narrowtable}

Resterna nerifrån och upp ger $13 = 1101_2$. Kontroll: $8 + 4 + 0 + 1 = 13$~\checkmark
\end{exempel}

\begin{exempel}[Från decimalt till hexadecimalt]
Omvandla $470$ till hexadecimal form. En rest mellan $10$ och $15$ skrivs med en bokstav.

\begin{narrowtable}{@{}ccc@{}}
  \thead{Division} & \thead{Kvot} & \thead{Rest} \\
  \midrule
  $470 / 16$ & $29$ & $6$ \\
  $29 / 16$ & $1$ & $13 = \mathrm{D}$ \\
  $1 / 16$ & $0$ & $1$ \\
\end{narrowtable}

Resterna nerifrån och upp ger $470 = \mathrm{1D6}_{16}$. Kontroll:
$1 \cdot 256 + 13 \cdot 16 + 6 = 470$~\checkmark
\end{exempel}

\begin{aside}
\note[Tips] Miniräknaren ger kvoten som decimaltal. $470 / 16 = 29{,}375$ betyder kvoten $29$ och
resten $0{,}375 \cdot 16 = 6$.
\end{aside}

\subsection{Mellan binärt och hexadecimalt}
\label{c5:sec:binarthex}
Eftersom $16 = 2^4$ motsvarar varje hexadecimal siffra exakt fyra bitar:

\begin{narrowtable}{@{}cc@{\hspace{2em}}cc@{\hspace{2em}}cc@{\hspace{2em}}cc@{}}
  \thead{Hex} & \thead{Binärt} & \thead{Hex} & \thead{Binärt} &
  \thead{Hex} & \thead{Binärt} & \thead{Hex} & \thead{Binärt} \\
  \midrule
  $0$ & $0000$ & $4$ & $0100$ & $8$ & $1000$ & $\mathrm{C}$ & $1100$ \\
  $1$ & $0001$ & $5$ & $0101$ & $9$ & $1001$ & $\mathrm{D}$ & $1101$ \\
  $2$ & $0010$ & $6$ & $0110$ & $\mathrm{A}$ & $1010$ & $\mathrm{E}$ & $1110$ \\
  $3$ & $0011$ & $7$ & $0111$ & $\mathrm{B}$ & $1011$ & $\mathrm{F}$ & $1111$ \\
\end{narrowtable}

\textbf{Från binärt till hexadecimalt:} dela in bitarna i grupper om fyra, \textbf{från höger},
fyll vid behov på med nollor till vänster och byt varje grupp mot sin hexadecimala siffra:
\[
  1011\,0110_2 = \underbrace{1011}_{\mathrm{B}}\;\underbrace{0110}_{6} = \mathrm{B6}_{16}
\]
\textbf{Från hexadecimalt till binärt:} byt varje siffra mot sina fyra bitar:
\[
  \mathrm{4D}_{16} = \underbrace{0100}_{4}\;\underbrace{1101}_{\mathrm{D}} = 0100\,1101_2
\]
En byte motsvarar alltså exakt två hexadecimala siffror, från \code{0x00} till \code{0xFF}.
Därför skrivs till exempel innehållet i en mikrokontrollers register ofta hexadecimalt:
\code{0xB6} är lättare att läsa än \code{0b10110110}.

\begin{aside}
\note[Obs] I datorsammanhang används ibland kilo för $2^{10} = 1\,024$ i stället för $1\,000$.
Den entydiga beteckningen är \textbf{kibi} (Ki): $1\,\text{KiB} = 1\,024$~byte.
\end{aside}
